% RASCUNHO -- IEEE CEC 2027 (conferencia). NAO COMPILADO (sem LaTeX na maquina em 05/10/2026).
% Todo numero deste texto tem rastro em RASTREIO_NUMEROS.md. Itens "% TODO" bloqueiam a submissao.
% Estudo 3 (conjunto estendido de operadores) esta em execucao: secao marcada como PENDENTE.
\documentclass[conference]{IEEEtran}
\usepackage{amsmath,amssymb}
\usepackage{booktabs}
\usepackage{graphicx}
\usepackage{cite}

\begin{document}

\title{Structure over Learning: A Pre-Specified Study of Adaptive Operator Selection in Variable Neighborhood Search for the CVRP}

% TODO(autores): nomes, afiliacoes, e-mails.
\author{\IEEEauthorblockN{[Author names pending]}
\IEEEauthorblockA{[Affiliations pending]}}

\maketitle

\begin{abstract}
Reinforcement-learning-guided metaheuristics are often credited with improving the search through adaptive operator selection (AOS), yet the learning component is rarely separated from the search structure it is embedded in. We reproduce a recent bandit-guided Variable Neighborhood Search for the Capacitated Vehicle Routing Problem (HG-UCB) and evaluate it in controlled, pre-specified studies on 74 CVRPLIB instances under equal operator-application budgets. First, the published reward, expressed in raw cost units, collapses the policy's softmax; normalizing it by the current solution cost is necessary: without normalization the method loses 1.0--1.3 percentage points (pp) of gap to the normalized version, on every instance. Second, with the reward corrected, HG-UCB is statistically indistinguishable from uniform random selection and from classical UCB1, $\varepsilon$-greedy and probability-matching rules. Third, a factorial decomposition shows that search structure matters more than learning: stopping the inner loop at a local optimum (0.5--0.6~pp) and randomizing the operator order (0.3--0.5~pp) each improve results more than any selection policy does, and no adaptive policy beats random order inside a standard VND. A third study extends the pool to ten operators of very unequal cost: the larger pool improves the gap by 0.5~pp at equal distance-evaluation budget, but adaptive selection---including a cost-aware reward and a patience-based stopping rule that lets it skip operators---still does not beat random selection, and a fixed order by cost is worse than random. A state-of-the-art reference solves the same instances to a 0.01--0.04\% gap, so all VNS variants studied here remain far from competitive; our claims concern the relative value of learning, not absolute performance.
\end{abstract}

\begin{IEEEkeywords}
adaptive operator selection, multi-armed bandits, variable neighborhood search, capacitated vehicle routing, reproducibility
\end{IEEEkeywords}

\section{Introduction}
Machine learning is increasingly embedded in metaheuristics to select search operators, tune parameters or guide moves~\cite{karimimamaghan2022}. Adaptive operator selection (AOS) based on multi-armed bandits has a long history in evolutionary computation~\cite{fialho2010,thierens2005} and in adaptive large neighborhood search~\cite{ropke2006}. A recent method couples a hybrid UCB1/gradient-bandit policy (HG-UCB) with a Variable Neighborhood Search (VNS)~\cite{hansen2001} for the Capacitated Vehicle Routing Problem (CVRP) and reports improvements over a plain VNS~\cite{kalatzantonakis2023}. A fair reading of such a result requires separating two things that are usually confounded: the \emph{selection policy} and the \emph{search structure} (stopping rule, acceptance rule, operator order) it operates in.

We ask: \emph{under an equal operator-application budget, does the bandit's learning contribute, and how large is that contribution relative to structural choices of the search?} We answer with three studies, each pre-specified in a dated research log before any data was produced.

\textbf{Contributions.}
(i)~A reproduction of HG-UCB-guided VNS that documents three interpretation gaps of the published pseudocode and identifies a reward-scale failure (Section~\ref{sec:repro}).
(ii)~A controlled comparison of HG-UCB against uniform selection and three classical AOS rules on 74 instances, two budgets and 10 seeds (Study~1).
(iii)~A factorial decomposition of the search structure and a test of AOS inside a standard VND (Study~2).
(iv)~A test with an extended, heterogeneous pool of ten operators, under a budget measured in distance evaluations, with a cost-aware reward and two stopping rules (Study~3).
(v)~An open protocol: pre-specification, exact regression controls, brute-force-validated operators and a state-of-the-art reference. % TODO(link): codigo e dados com DOI persistente.

\section{Background}
\subsection{Problem and benchmark}
The CVRP asks for a set of routes, each starting and ending at a depot and serving customers with demands, such that each customer is served once, no route exceeds the vehicle capacity, and the total distance is minimized. We use 74 instances from CVRPLIB: sets A (26) and B (23) with an exact number of routes, set CMT (4), and the X set (21)~\cite{uchoa2017}. % TODO(citacoes): Augerat et al. 1995 e Christofides et al. 1979 pendentes de verificacao (ver refs.bib).

\subsection{HG-UCB-guided VNS}
Five neighborhoods are used: 2-Opt and Move(1) (intra-route) and Swap(1,1), Relocate(1,0) and Relocate(2,0) (inter-route)~\cite{kalatzantonakis2023}. At each step the VNS perturbs the incumbent (shake) and an inner loop repeatedly selects a neighborhood, applies the best improving move and updates the bandit. HG-UCB selects with a UCB1 rule~\cite{auer2002} whose exploitation term is a softmax over gradient-bandit preferences, with a cold start, an $\varepsilon$-greedy wrapper and a drift detector that resets the policy.

\section{Reproduction and a reward-scale failure}\label{sec:repro}
Our plain GVNS baseline reaches a mean gap of 3.35\% on sets A, B and X, against 3.39\% reported for the corresponding column of the original study. Reproducing the bandit variant required resolving three gaps in the pseudocode: the inner-loop stopping condition $\max(q)<b_t$ is unreachable with five operators and $b_t=0.2$ (a softmax of five values never falls below $1/5$), so an explicit cap on inner iterations is needed; the scope of the reference solution $x'$ within a step is ambiguous; and the text and the pseudocode disagree on what feeds the drift detector (Q-values or the raw reward); we use the raw reward, because feeding Q-values worsened the gap in our tests. We adopted a cap of 40 inner iterations and a progressive $x'$ (Section~\ref{sec:protocol}).

With the reward of the original formulation, expressed in absolute cost units, and the published learning rate $\alpha=5.0$, the softmax over preferences saturates to a nearly one-hot distribution in 99.6\% of the iterations of an 8-instance diagnostic sample, neutralizing both the probability filter and the exploration mechanism. Dividing the positive reward by the current solution cost makes it dimensionless and reduces saturation to 12.5\% at the same $\alpha$. We use this normalized reward (with $\alpha=0.5$, chosen on the same instances; see Section~\ref{sec:threats}) as ``HG-UCB'' hereafter, and the unnormalized, $\alpha=5.0$ variant as ``HG-UCB literal''.

\section{Experimental protocol}\label{sec:protocol}
\textbf{Budget.} The budget is the number of operator applications ($B\in\{800,4000\}$ in Studies 1--2), counted by a wrapper around each neighborhood's best-improvement function; the shake does not count. The production configuration (20 steps $\times$ 40 inner iterations) uses exactly 800 applications. \textbf{Runs.} 10 seeds per instance, the same Clarke--Wright initial solution for all arms. \textbf{Metric.} Gap (\%) to the reference cost, averaged over seeds per instance; the unit of pairing is the instance ($n=74$). \textbf{Inference.} Paired two-sided Wilcoxon signed-rank tests, 95\% bootstrap confidence intervals over instances (10\,000 resamples), Holm correction over a family fixed before the runs. Effects are reported as paired mean differences in percentage points (negative favors the first arm). \textbf{Pre-specification.} Hypotheses, arms, tests and family size were written in a dated research log before each study ran; no external registry timestamp exists. % TODO(honestidade): considerar deposito em OSF/Zenodo e git para dar timestamp externo.
\textbf{Controls.} The generic search loop reproduces the two earlier search structures exactly (identical costs on all compared runs), operator counts equal $B$ in every run, and all solutions are feasible.

\section{Study 1: does the bandit beat chance?}
Within the original loop (fixed 40 inner applications per step), we compare HG-UCB with uniform selection, UCB1, $\varepsilon$-greedy, probability matching and classical GVNS (fixed-order VND with a matched application budget). Table~\ref{tab:s1} reports mean gaps. HG-UCB is not better than uniform selection at either budget ($+0.06$~pp, adjusted $p=1.0$ at $B=800$; $+0.10$~pp, 95\% CI $[0.01,0.19]$, adjusted $p=0.11$ at $B=4000$; Holm over 12 tests), nor does it differ significantly from UCB1, $\varepsilon$-greedy or probability matching (adjusted $p\ge 0.15$). The literal variant is $0.96$~pp ($B=800$) and $1.31$~pp ($B=4000$) worse than normalized HG-UCB, losing on all 74 instances at both budgets. Classical GVNS has the lowest mean gap of all arms and is significantly better than HG-UCB ($0.46$--$0.54$~pp; better on 63 of 74 instances at both budgets), a difference examined in Study~2.

\begin{table}[t]
\caption{Study 1: mean gap (\%) over 74 instances, original loop.}\label{tab:s1}
\centering
\begin{tabular}{lcc}
\toprule
Policy & $B=800$ & $B=4000$ \\
\midrule
Classical GVNS & 3.01 & 2.30 \\
Uniform & 3.42 & 2.74 \\
UCB1 & 3.42 & 2.75 \\
Prob.\ matching & 3.47 & 2.79 \\
HG-UCB (normalized) & 3.47 & 2.84 \\
$\varepsilon$-greedy & 3.50 & 2.87 \\
HG-UCB literal & 4.43 & 4.15 \\
\bottomrule
\end{tabular}
\end{table}

A subset experiment at $B=20\,000$ (22 stratified instances, 5 seeds, three policies; exploratory, low power) shows the same pattern: HG-UCB minus uniform is $-0.06$, $-0.11$ and $+0.21$~pp at $B=800$, $4000$ and $20\,000$, all with confidence intervals that include zero, and classical GVNS remains about $0.5$~pp better than HG-UCB.

\section{Study 2: where does the difference come from?}
A diagnostic shows why the original loop is inefficient: 85--90\% of the 40 inner applications per step occur after the last improvement of that step (3 instances, 1 seed, exploratory); since every neighborhood is deterministic, repeating them at a local optimum cannot change the solution. Classical GVNS performs 60--165 shakes within 800 applications, the original loop only 20. We therefore decompose the structure with three factors: operator \emph{order} (fixed or uniformly random), \emph{acceptance} (accept the local-search result if it improves the best solution, ``best'', or only the shaken point, ``shaken'') and inner-loop \emph{termination} (stop when all five operators failed since the last improvement, ``LO'', or run 40 applications, ``F40'').

\begin{table}[t]
\caption{Study 2A: mean gap (\%) over 74 instances by structure.}\label{tab:s2a}
\centering
\begin{tabular}{lccc}
\toprule
Order / acceptance / termination & $B=800$ & $B=4000$ \\
\midrule
Random / best / LO & 2.69 & 1.94 \\
Random / shaken / LO & 2.90 & 2.11 \\
Fixed / best / LO (classical GVNS) & 3.01 & 2.30 \\
Fixed / shaken / LO & 3.21 & 2.60 \\
Random / best / F40 & 3.33 & 2.54 \\
Random / shaken / F40 (original loop) & 3.42 & 2.74 \\
\bottomrule
\end{tabular}
\end{table}

Termination is the dominant factor: LO versus F40 improves the gap by $0.64$ and $0.52$~pp at $B=800$ and by $0.60$ and $0.62$~pp at $B=4000$ (two acceptance rules; better on 66--70 of 74 instances; Holm-adjusted $p<10^{-4}$). Acceptance contributes $0.09$--$0.21$~pp (the LO contrast at $B=4000$ does not survive correction, adjusted $p=0.054$). Contrary to our pre-specified expectation, random order is better than fixed order by $0.33$--$0.36$~pp ("best") and $0.31$--$0.49$~pp ("shaken"), significantly at both budgets: the fixed order always restarts from 2-Opt, which then absorbs 39\% of the applications. Moving from the original loop to random order, ``best'' acceptance and LO termination reduces the gap from 2.74\% to 1.94\% at $B=4000$ without any learning.

\textbf{AOS inside VND (2B).} In the best structure (``best''/LO, selection among not-yet-failed operators) we compare four adaptive policies against random order. UCB1 and probability matching tie with random order (differences $\le 0.09$~pp, not significant); $\varepsilon$-greedy ($+0.23$ and $+0.19$~pp) and HG-UCB ($+0.16$ and $+0.14$~pp) are significantly \emph{worse} at both budgets. Adaptive policies appear better than the classical fixed order (by $0.10$--$0.31$~pp), but this reflects the order effect above, not adaptation. Of 28 pre-specified tests, 21 are significant after Holm correction.

\section{Study 3: an extended, heterogeneous operator pool}
The five operators are similar and deterministic, which could leave no room for adaptation. We therefore add five operators of different nature and cost---2-opt*, Or-opt (segments of 2--3), intra-route swap, CROSS exchange (segments up to 3) and SWAP*~\cite{vidal2022} (a naive implementation)---for a pool of ten, each validated against an independent brute-force enumeration (1\,600 tests). Their cost differs by almost two orders of magnitude (about 0.6k to 35k distance evaluations per application at $n=100$) and, at local optima of the original five operators, only SWAP* found improvements in our profiling sample. A budget in operator applications would therefore treat SWAP* as free; we measure the budget in \emph{distance evaluations} instead (a deterministic proxy for time, with roughly constant cost per evaluation): $E_1=600\,n^2$ and $E_2=5E_1$, where $E_1$ approximates the effort of 800 applications of the original pool. As a consistency check, the original pool with random order and LO termination reaches 1.98\% at $E_2$ in this study, against 1.94\% at $B=4000$ in Study~2. The shake uses only the original five operators, so the perturbation is unchanged.

We cross two stopping rules---LO (all operators must fail, so adaptation can only reorder) and R5 (stop after five consecutive failures, sampling with replacement, so adaptation can avoid costly useless operators)---with five selection policies: a fixed order by increasing cost (FOC, fixed before the runs), uniform, and UCB1, probability matching and HG-UCB with a cost-aware reward (the Study~1--2 reward multiplied by the ratio between the running mean cost per application and the application's cost, clipped to $[0.1,10]$). Two references use the original pool under the same budget. Of 34 pre-specified tests (Holm), 13 are significant (3 at $E_1$, 10 at $E_2$).

\begin{table}[t]
\caption{Study 3: mean gap (\%) over 74 instances (budgets $E_1$, $E_2$).}\label{tab:s3}
\centering
\begin{tabular}{lcc}
\toprule
Arm & $E_1$ & $E_2$ \
\midrule
R5, UCB1 (cost-aware) & 2.13 & 1.34 \
R5, uniform & 2.22 & 1.37 \
R5, prob.\ matching & 2.22 & 1.42 \
LO, uniform & 2.26 & 1.45 \
LO, prob.\ matching & 2.27 & 1.45 \
LO, UCB1 & 2.29 & 1.46 \
LO, HG-UCB & 2.36 & 1.56 \
R5, HG-UCB & 2.18 & 1.57 \
LO, fixed order by cost & 2.40 & 1.64 \
R5, fixed order by cost & 2.18 & 1.65 \
\midrule
Original pool, LO, uniform & 2.76 & 1.98 \
Original pool, LO, fixed order & 3.02 & 2.32 \
\bottomrule
\end{tabular}
\end{table}

The larger pool is the only large effect: with the same stopping rule and random order, ten operators beat five by $0.50$ and $0.53$~pp at $E_1$ and $E_2$ (better on 59 and 58 of 74 instances; adjusted $p<10^{-4}$), in all four instance sets. Random order again beats the classical fixed order in the original pool ($0.26$ and $0.34$~pp). A fixed order by cost is \emph{worse} than random order with the large pool ($+0.13$ and $+0.19$~pp under LO; $+0.28$~pp under R5 at $E_2$), plausibly because the cheap operators rarely improve while the productive one comes last. No adaptive policy beats uniform selection: UCB1 and probability matching tie with it under both stopping rules and budgets, and HG-UCB is significantly worse at $E_2$ ($+0.11$~pp under LO, $+0.20$~pp under R5). UCB1 and probability matching do beat the fixed order by cost at $E_2$ ($0.18$--$0.31$~pp), which reflects the weakness of that order, not adaptation. The cost-aware reward made no measurable difference to HG-UCB ($-0.002$ and $-0.042$~pp, not significant). The uniform and adaptive policies spend 47--61\% of their distance evaluations in SWAP* (the fixed order by cost, 24--45\%), with HG-UCB concentrating most (61\% under LO), consistent with SWAP* being the only operator that found improvements at local optima in our profiling sample. The stopping rule (R5 versus LO) has no detectable effect under random order (differences of $0.04$--$0.08$~pp, not significant).

\section{Reference solver and context}
As a ceiling, HGS-CVRP~\cite{vidal2022}, run for the same wall-clock budget as our VNS (0.2--10~s per instance) reaches a mean gap of 0.035\%, and 0.010\% with 0.1~s per customer, against 1.9--3.5\% for the (non-literal) VNS variants here. Part of this distance is implementation language (C++ versus Python), so we do not claim a ranking, only that no variant studied is competitive. An attention-based constructive model~\cite{kool2019} used zero-shot reaches 9.43\% (58 of 74 instances feasible under the exact-route criterion).

\section{Threats to validity and limitations}\label{sec:threats}
(1) Study~3 uses a pool of ten operators, still deterministic best-improvement moves from one family of neighborhoods; richer or stochastic operators, contextual policies or different reward definitions may behave differently. (2) The learning rate $\alpha=0.5$ was selected on the same 74 instances, which favors HG-UCB, while classical policies use untuned literature defaults, and the patience ($r=5$), the cost-aware reward and its clipping were fixed without tuning; a tuned classical policy might differ. (3) The budget counts operator applications, which ignores cost differences between operators (roughly 8--12\,ms per application on average across arms here). (4) All variants are far from state-of-the-art quality on a Python platform. (5) Wilcoxon tests on instance means ignore within-instance seed variance. (6) Pre-specification lives in a dated research log without an external timestamp. (7) Results are for CVRPLIB instances only.

\section{Conclusion}
For a bandit-guided VNS on the CVRP, the learning component contributed nothing measurable once its reward was normalized, while two structural choices of the search---stopping at local optima and randomizing operator order---improved the gap by 0.8~pp in combination. Extending the pool to ten operators of unequal cost improved the gap by a further 0.5~pp, yet no adaptive policy beat uniform selection there either, even with a cost-aware reward and a stopping rule that allows operators to be skipped. We therefore recommend that studies of learned operator selection report a matched non-adaptive control, the search structure's own ablation and the operator pool's contribution.

% TODO(fechamento): declaracao de uso de IA generativa (ferramenta, finalidade, supervisao humana);
% disponibilidade de dados/codigo com link persistente; conflito de interesse; financiamento.

\bibliographystyle{IEEEtran}
\bibliography{refs}

\end{document}
